\clearpage
\thispagestyle{empty}

\vspace*{0.8in}

This dissertation is concerned with the design of computing architectures which
support the implementation of \emph{declarative} programming languages.
Declarative languages, which include the functional and logic programming
languages, are based on mathematical theories of computation. These mathematical
foundations endow declarative languages with elegant semantics and
expressiveness, but make efficient implementation on conventional von-Neumann
computers difficult.

An alternative computing architecture, \RED{2}, is developed from first
principles to support the efficient implementation of functional programming
languages by directly executing their underlying mathematical theory, the
lambda calculus. The machine language of \RED{2} is THOR, a high-level language
which supports partial evaluation and incremental reduction. Two novel features
of the \RED{2} architecture are a program counter which moves both forwards and
backwards and an incremental memory management policy which eliminates the need
for explicit garbage collection. The strengths and weaknesses of \RED{2} are
discussed in detail, and several enhancements are presented.

Finally, the possibility of extending \RED{2} to incorporate Horn clause logic
programming is explored.
